\section{Synthetic recovery tests}\label{sec:synthetic}

Before applying the fitter to real data, we test it on synthetic CLASS-like spectra.

\paragraph{Generator} The synthetic spectra are generated with the full forward model, including secondary fluorescence, which the fitted model omits. The calibration is the December 2024 one, the background uses the real December 2024 particle and instrumental templates ($\beta=1$, $\gamma=0.5$), and each spectrum is an \SI{8}{s}, 16-SCD exposure with Poisson noise. Compositions are drawn uniformly from lunar ranges (Mg/Si 0.10--0.35, Al/Si 0.20--0.80, Ca/Si 0.30--0.60), with incidence angles of 10--80\textdegree. The scattered-solar amplitude is set to match the continuum level of the December 2024 stack.

\paragraph{Results at the correct temperature} The signal level is expressed as the number of Si K$\alpha$ counts. For a \SI{5}{MK} solar spectrum (the median GOES temperature of 2024--2026), \cref{tab:synthetic} lists the median relative bias and robust (MAD-based) relative scatter of the recovered ratios. It also lists the width of the pull distribution, $(\hat r-r)/\hat\sigma_r$. Mg/Si and Al/Si are recovered with $|\text{bias}|\le\SynMaxBiasMgAl\%$ for every signal level, and with pulls of width \SynPullLo--\SynPullHi, so the Fisher uncertainties are realistic. Omitting secondary fluorescence and fitting the instrumental amplitude freely do not bias them. Ca/Si is biased by +30\% even when the temperature is exact, because a \SI{5}{MK} spectrum barely excites Ca.

\paragraph{Temperature error} The main systematic is the solar temperature. We generate spectra at one temperature and fit them assuming another, to mimic error in the GOES temperature. With a true spectrum 30\% hotter than assumed, Mg/Si is biased by \SynHotMg\% and Al/Si by \SynHotAl\%. With a true spectrum 20\% cooler, the biases are \SynColdMg\% and \SynColdAl\%. Ca/Si changes by a factor of $\sim$9 or collapses to zero. These numbers set the systematic uncertainty quoted for the maps (\cref{sec:limitations}). They also show why Ca/Si, Ti/Si and Fe/Si cannot be reported without a measured solar spectrum.

\begin{table}[t]
\centering
\caption{Synthetic recovery at a \SI{5}{MK} solar temperature (80 spectra per level): median relative bias / robust relative scatter (\%), and pull width.}\label{tab:synthetic}
\begin{tabular}{lrrrr}
\toprule
Si K$\alpha$ counts & 300 & 1,000 & 3,000 & 10,000\\
\midrule
Mg/Si bias / scatter (\%) & +2 / 21 & $-$3 / 12 & $-$1 / 8 & $-$1 / 4\\
Mg/Si pull width & 0.78 & 1.03 & 1.07 & 1.21\\
Al/Si bias / scatter (\%) & +1 / 39 & 0 / 17 & 0 / 6 & 0 / 3\\
Al/Si pull width & 1.18 & 1.09 & 0.90 & 0.93\\
Ca/Si bias / scatter (\%) & +5 / 156 & +22 / 127 & +51 / 75 & +28 / 52\\
Ca/Si pull width & 0.69 & 0.98 & 1.12 & 1.67\\
Quality-0 fraction & 0.99 & 1.00 & 1.00 & 1.00\\
\bottomrule
\end{tabular}

\end{table}
